\chapter{Sterowalność, obserwowalność i stabilizowalność}
\label{ch:lin25-control}
\index{sterowalność@Sterowalność}
\index{obserwowalność@Obserwowalność}
\index{kryterium Kalmana@Kryterium Kalmana}
\index{kryterium PBH@Kryterium PBH}
\index{stabilizowalność@Stabilizowalność}
\index{wykrywalność@Wykrywalność}

\noindent\textbf{Cel rozdziału.}
W Rozdziale~\ref{ch:lin24-riccati}
dowiedliśmy stabilności optymalnych sprzężeń
przy założeniu, że macierz otwartej pętli
jest już stabilna. Chcemy wiedzieć,
kiedy sterowanie może ustabilizować
układ początkowo niestabilny
i kiedy jego stan można odtworzyć
z obserwowanego wyjścia.
Prowadzi to do sterowalności,
obserwowalności, ich kryteriów
macierzowych i rozkładu Kalmana.
Udowodnimy również spektralne
kryteria stabilizowalności
i wykrywalności, nie zakładając
bez dowodu twierdzeń o rozmieszczaniu
biegunów sprzężenia zwrotnego.

Pracujemy nad ciałem \(\mathbb C\);
rzeczywiste odpowiedniki zostaną
wskazane w osobnej uwadze.
Ustalamy liczby całkowite
\(n,m,p\ge1\) oraz macierze
\(A\in\operatorname{Mat}_{n\times n}(\mathbb C)\),
\(B\in\operatorname{Mat}_{n\times m}(\mathbb C)\),
\(C\in\operatorname{Mat}_{p\times n}(\mathbb C)\).
Stan to \(x\in\mathbb C^n\),
sterowanie to \(u\in\mathbb C^m\),
a obserwacja to \(y=Cx\in\mathbb C^p\).
Przez \(\operatorname{im}T\)
i \(\ker T\) oznaczamy odpowiednio
obraz i jądro odwzorowania liniowego \(T\).

\section{Sterowalność w czasie dyskretnym}
\label{sec:lin25-reach}

\begin{definicja}[Macierz i przestrzeń sterowalności]
\label{def:lin25-kalman}
Dla pary \((A,B)\) definiujemy
\emph{macierz sterowalności Kalmana}
oraz \emph{przestrzeń osiągalną}
wzorami
\[
 \mathcal K(A,B)=
 \begin{bmatrix}
 B&AB&A^2B&\cdots&A^{n-1}B
 \end{bmatrix},
 \qquad
 \mathcal R(A,B)=\operatorname{im}\mathcal K(A,B).
\]
Parę \((A,B)\) nazywamy
\emph{sterowalną}, jeśli
\(\mathcal R(A,B)=\mathbb C^n\),
równoważnie
\(\operatorname{rank}\mathcal K(A,B)=n\).
\end{definicja}

\begin{lemat}[Niezmienniczość przestrzeni osiągalnej]
\label{lem:lin25-invariant}
Dla każdej pary \((A,B)\)
przestrzeń \(\mathcal R(A,B)\)
zawiera \(\operatorname{im}B\),
jest niezmiennicza względem \(A\)
i spełnia
\[
 \mathcal R(A,B)
 =\sum_{k=0}^{\infty}\operatorname{im}A^kB.
\]
Jest ona najmniejszą
\(A\)-niezmienniczą podprzestrzenią
zawierającą \(\operatorname{im}B\).
\end{lemat}
\begin{proof}
Z definicji
\(\operatorname{im}B\subseteq\mathcal R(A,B)\).
Dla \(0\le k<n-1\)
mamy
\(A(\operatorname{im}A^kB)
 =\operatorname{im}A^{k+1}B
 \subseteq\mathcal R(A,B)\).
Pozostaje sprawdzić \(k=n-1\).
Z twierdzenia Cayleya--Hamiltona,
udowodnionego w
Rozdziale~\ref{ch:lin16-niezmiennicze},
istnieją współczynniki
\(c_0,\ldots,c_{n-1}\in\mathbb C\)
takie, że
\[
 A^n=-\sum_{j=0}^{n-1}c_jA^j.
\]
Po pomnożeniu przez \(B\)
otrzymujemy
\(\operatorname{im}A^nB
 \subseteq\mathcal R(A,B)\).
Wobec tego \(\mathcal R(A,B)\)
jest \(A\)-niezmiennicza,
a więc zawiera obrazy
\(A^kB\) dla wszystkich \(k\ge0\).
Odwrotne zawieranie
wynika bezpośrednio z definicji.
Każda \(A\)-niezmiennicza
podprzestrzeń zawierająca
\(\operatorname{im}B\)
musi zawierać także wszystkie
\(\operatorname{im}A^kB\);
stąd minimalność.
\end{proof}

\begin{twierdzenie}[Kryterium Kalmana w czasie dyskretnym]
\label{thm:lin25-kalman}
Niech \(n,m\ge1\) będą liczbami całkowitymi,
\(A\in\operatorname{Mat}_{n\times n}(\mathbb C)\)
oraz \(B\in\operatorname{Mat}_{n\times m}(\mathbb C)\).
Rozważmy dla \(k\ge0\) układ
\(x_{k+1}=Ax_k+Bu_k\),
gdzie \(x_k\in\mathbb C^n\)
i \(u_k\in\mathbb C^m\).
Dla ustalonego \(N\ge1\)
zbiór stanów osiągalnych
w chwili \(N\) ze stanu \(x_0=0\)
jest równy
\[
 \operatorname{im}
 \begin{bmatrix}
 B&AB&\cdots&A^{N-1}B
 \end{bmatrix}.
\]
Następujące warunki
są równoważne:
\begin{enumerate}[label=\textup{(\roman*)}]
 \item para \((A,B)\) jest sterowalna;
 \item każdy \(z\in\mathbb C^n\)
 można osiągnąć ze stanu \(0\)
 w dokładnie \(n\) krokach;
 \item istnieje \(N\ge1\),
 dla którego każdy
 \(z\in\mathbb C^n\)
 można osiągnąć z \(0\)
 w dokładnie \(N\) krokach;
 \item dla każdych \(x,z\in\mathbb C^n\)
 istnieje ciąg sterowań,
 który przeprowadza \(x\)
 w \(z\) w dokładnie \(n\) krokach.
\end{enumerate}
\end{twierdzenie}
\begin{proof}
Iterując równanie stanu
otrzymujemy kolejno
\(x_1=Ax_0+Bu_0\),
\(x_2=A^2x_0+ABu_0+Bu_1\).
Indukcja daje dla każdego \(N\ge1\)
\[
 x_N=A^Nx_0+
      \sum_{j=0}^{N-1}A^{N-1-j}Bu_j.
\]
Po podstawieniu \(x_0=0\)
i odwróceniu kolejności
wektorów \(u_j\)
otrzymujemy dokładnie obraz
macierzy podanej w tezie.
Dla \(N=n\) jest to
\(\mathcal R(A,B)\),
więc (i) i (ii) są równoważne,
a (ii) implikuje (iii).

Jeśli zachodzi (iii),
to obraz wskazanej macierzy
jest całą \(\mathbb C^n\);
na mocy
Lematu~\ref{lem:lin25-invariant}
jest jednak zawarty
w \(\mathcal R(A,B)\).
Zatem (i).

Załóżmy (i).
Dla dowolnych \(x,z\)
połóżmy \(w=z-A^nx\).
Z (ii) można w \(n\) krokach
otrzymać ze stanu zerowego
dowolny wektor \(w\).
Wybierając te same sterowania
dla stanu początkowego \(x\),
otrzymujemy
\(x_n=A^nx+w=z\).
To daje (iv).
Z (iv) wynika (ii)
po przyjęciu \(x=0\).
\end{proof}

\begin{przyklad}[Łańcuch dwóch stanów]
\label{prz:lin25-chain}
Dla
\[
 A=\begin{pmatrix}0&1\\0&0\end{pmatrix},
 \qquad
 B=\begin{pmatrix}0\\1\end{pmatrix}
\]
mamy
\(AB=(1,0)^{\mathsf T}\),
więc
\[
 \mathcal K(A,B)
 =\begin{pmatrix}0&1\\1&0\end{pmatrix},
 \qquad
 \det\mathcal K(A,B)=-1.
\]
Para jest sterowalna.
Zerowy stan przechodzi
po dwóch krokach w
\[
 x_2=ABu_0+Bu_1
     =\begin{pmatrix}u_0\\u_1\end{pmatrix}.
\]
Aby osiągnąć \(z=(a,b)^{\mathsf T}\),
wystarczy więc przyjąć
\(u_0=a\), \(u_1=b\).
W pierwszym kroku,
gdy \(x_0=0\),
osiągalne są jedynie
wektory postaci \((0,b)^{\mathsf T}\);
nie należy zatem zastępować
warunku „w \(n\) krokach”
warunkiem „w jednym kroku”.
\end{przyklad}

\section{Sterowalność w czasie ciągłym}
\label{sec:lin25-reach-cont}

\begin{definicja}[Ciągły układ sterowany]
\label{def:lin25-cont-system}
Niech \(T>0\).
Rozpatrujemy równanie
\[
 x'(t)=Ax(t)+Bu(t),\qquad
 x(0)=x_0,\qquad 0\le t\le T,
\]
dla przedziałami ciągłego
sterowania
\(u\colon[0,T]\to\mathbb C^m\).
Parę \((A,B)\)
nazywamy \emph{sterowalną
w czasie ciągłym}, jeśli
dla każdej liczby \(T>0\)
oraz dowolnych \(x_0,z\in\mathbb C^n\)
istnieje takie sterowanie,
że rozwiązanie spełnia \(x(T)=z\).
\end{definicja}

\begin{lemat}[Wzór na rozwiązanie z wymuszeniem]
\label{lem:lin25-variation}
Dla macierzy \(A,B\)
i przedziałami ciągłego sterowania
\(u\colon[0,T]\to\mathbb C^m\)
jedyne ciągłe,
przedziałami różniczkowalne
rozwiązanie równania
z Definicji~\ref{def:lin25-cont-system}
ma postać
\[
 x(t)=e^{tA}x_0+
      \int_0^t e^{(t-s)A}Bu(s)
      \,\mathrm ds,\qquad 0\le t\le T.
\]
\end{lemat}
\begin{proof}
Jeśli \(x\) jest rozwiązaniem,
pomnóżmy równanie
przez \(e^{-tA}\).
Z reguły iloczynu
i równości
\(\frac{\mathrm d}{\mathrm dt}e^{-tA}
 =-Ae^{-tA}\),
udowodnionej w
Rozdziale~\ref{ch:lin18-jordan-chevalley},
otrzymujemy na każdym przedziale
ciągłości \(u\)
\[
 \frac{\mathrm d}{\mathrm dt}
    (e^{-tA}x(t))
 =-Ae^{-tA}x(t)+e^{-tA}(Ax(t)+Bu(t))
 =e^{-tA}Bu(t).
\]
Całkując kawałkami od \(0\) do \(t\)
i mnożąc przez \(e^{tA}\),
otrzymujemy wzór z tezy.
Odwrotnie, różniczkując
podany wzór na przedziałach,
gdzie \(u\) jest ciągłe,
otrzymujemy równanie stanu
i \(x(0)=x_0\).
Dla dwóch rozwiązań
ich różnica \(v\)
spełnia \(v'=Av\) i \(v(0)=0\),
więc \(v(t)=e^{tA}v(0)=0\).
To zapewnia jednoznaczność.
\end{proof}

\begin{definicja}[Gramian sterowalności]
\label{def:lin25-gramian}
Dla każdej \(T>0\)
\emph{gramianem sterowalności}
nazywamy macierz hermitowską
\[
 W_T=\int_0^T
      e^{tA}BB^*e^{tA^*}
      \,\mathrm dt.
\]
\end{definicja}

\begin{twierdzenie}[Gramian i kryterium sterowalności ciągłej]
\label{thm:lin25-gramian}
Dla każdej pary macierzy
\(A\in\operatorname{Mat}_{n\times n}(\mathbb C)\),
\(B\in\operatorname{Mat}_{n\times m}(\mathbb C)\)
następujące warunki
są równoważne:
\begin{enumerate}[label=\textup{(\roman*)}]
 \item \(\operatorname{rank}\mathcal K(A,B)=n\);
 \item dla każdego \(T>0\)
 gramian \(W_T\) jest dodatnio
 określony;
 \item istnieje \(T>0\),
 dla którego \(W_T\) jest
 dodatnio określony;
 \item para \((A,B)\)
 jest sterowalna w czasie ciągłym.
\end{enumerate}
Jeśli \(W_T>0\), to
dla dowolnych \(x_0,z\in\mathbb C^n\)
sterowanie
\[
 u(s)=B^*e^{(T-s)A^*}W_T^{-1}
             (z-e^{TA}x_0),\quad0\le s\le T,
\]
przenosi stan \(x_0\)
w stan \(z\) w chwili \(T\).
\end{twierdzenie}
\begin{proof}
Dla każdego \(v\in\mathbb C^n\)
mamy
\[
 v^*W_Tv
 =\int_0^T
      \bigl\|B^*e^{tA^*}v\bigr\|_2^2
      \,\mathrm dt\ge0.
\]
Ciągłość funkcji podcałkowej
pokazuje, że \(v^*W_Tv=0\)
dokładnie wtedy, gdy
\(B^*e^{tA^*}v=0\)
dla wszystkich \(t\in[0,T]\).

Jeżeli ta funkcja znika,
to po zróżniczkowaniu
jej szeregu wykładniczego
w punkcie \(t=0\) mamy
\(B^*(A^*)^kv=0\)
dla wszystkich całkowitych \(k\ge0\).
W szczególności dla
\(k=0,\ldots,n-1\)
otrzymujemy
\(\mathcal K(A,B)^*v=0\).
Odwrotnie, jeśli
\(\mathcal K(A,B)^*v=0\),
to \(B^*(A^*)^kv=0\)
dla \(0\le k<n\).
Z twierdzenia Cayleya--Hamiltona,
zastosowanego do \(A^*\),
wynika indukcyjnie
zerowanie także dla wszystkich
\(k\ge n\).
Szereg
\(B^*e^{tA^*}v
 =\sum_{k=0}^\infty
      \frac{t^k}{k!}B^*(A^*)^kv\)
jest więc tożsamościowo zerowy.
Dowiedliśmy, że
\[
 \ker W_T=\ker\mathcal K(A,B)^*.
\]
Równość obowiązuje dla
każdego \(T>0\).
Ponieważ \(W_T\ge0\),
dodatnia określoność
jest równoważna
\(\ker W_T=\{0\}\).
Stąd równoważność
(i), (ii) i (iii).

Załóżmy (ii) i wybierzmy
dowolne \(x_0,z,T\).
Niech \(d=z-e^{TA}x_0\)
i zdefiniujmy \(u\)
wzorem z tezy.
Z
Lematu~\ref{lem:lin25-variation}
otrzymujemy
\[
 \begin{aligned}
 x(T)-e^{TA}x_0
 &=\int_0^T e^{(T-s)A}BB^*
        e^{(T-s)A^*}W_T^{-1}d
       \,\mathrm ds\\
 &=W_TW_T^{-1}d=d,
 \end{aligned}
\]
bo podstawienie \(t=T-s\)
zmienia całkę w \(W_T\).
Zatem \(x(T)=z\);
otrzymaliśmy (iv).

Załóżmy (iv).
Dla dowolnego \(T>0\)
i każdego \(d\in\mathbb C^n\)
można ze stanu \(0\) osiągnąć
\(d\) w chwili \(T\).
Gdyby \(W_T\) nie był
dodatnio określony,
to z jego nieujemnej
określoności istniałby
\(v\ne0\) taki, że
\(B^*e^{tA^*}v=0\)
dla wszystkich \(0\le t\le T\).
Wówczas dla każdej trajektorii
startującej od zera
\[
 v^*x(T)=\int_0^T
          v^*e^{(T-s)A}Bu(s)
          \,\mathrm ds=0.
\]
Nie można by osiągnąć
na przykład \(d=v\),
gdyż \(v^*v>0\).
Sprzeczność dowodzi (ii).
\end{proof}

\begin{przyklad}[Gramian ciągły w wymiarze dwa]
\label{prz:lin25-gramian}
Niech \(A,B\) będą macierzami
z Przykładu~\ref{prz:lin25-chain}.
Ponieważ \(A^2=0\),
mamy \(e^{tA}=I+tA\),
a więc
\(e^{tA}B=(t,1)^{\mathsf T}\).
Dla dowolnego \(T>0\)
otrzymujemy
\[
 \begin{aligned}
 W_T&=\int_0^T
 \begin{pmatrix}t^2&t\\t&1\end{pmatrix}
 \,\mathrm dt\\
 &=\begin{pmatrix}
       T^3/3&T^2/2\\
       T^2/2&T
   \end{pmatrix}.
 \end{aligned}
\]
Macierz jest dodatnio określona:
jej lewy górny element
jest dodatni, a wyznacznik
wynosi
\[
 \det W_T=\frac{T^4}{3}
                -\frac{T^4}{4}
          =\frac{T^4}{12}>0.
\]
Stąd sterowalność obowiązuje
dla każdego, również bardzo
małego, dodatniego czasu \(T\).
Nie oznacza to,
że koszt energii potrzebnej
do sterowania pozostaje
ograniczony, gdy \(T\to0\):
wzór na \(W_T^{-1}\)
zawiera rosnące potęgi \(T^{-1}\).
\end{przyklad}

\section{Obserwowalność i jej związek ze sterowalnością}
\label{sec:lin25-observe}

\begin{definicja}[Macierz obserwowalności]
\label{def:lin25-observable}
Dla pary \((C,A)\)
definiujemy \emph{macierz
obserwowalności Kalmana}
oraz \emph{podprzestrzeń
nieobserwowalną}:
\[
 \mathcal O(C,A)=
 \begin{bmatrix}
 C\\CA\\CA^2\\\vdots\\CA^{n-1}
 \end{bmatrix},
 \qquad
 \mathcal N(C,A)=\ker\mathcal O(C,A).
\]
Parę \((C,A)\)
nazywamy \emph{obserwowalną},
jeśli \(\mathcal N(C,A)=\{0\}\),
równoważnie
\(\operatorname{rank}\mathcal O(C,A)=n\).
\end{definicja}

\begin{twierdzenie}[Kryterium obserwowalności dyskretnej]
\label{thm:lin25-observability}
Niech \(x_{k+1}=Ax_k\)
i \(y_k=Cx_k\)
dla wszystkich \(k\ge0\).
Następujące warunki
są równoważne:
\begin{enumerate}[label=\textup{(\roman*)}]
 \item para \((C,A)\) jest obserwowalna;
 \item wektor \(x_0\)
 jest jednoznacznie wyznaczony
 przez obserwacje
 \(y_0,y_1,\ldots,y_{n-1}\);
 \item jeżeli wszystkie
 \(y_0,\ldots,y_{n-1}\)
 są zerowe, to \(x_0=0\);
 \item jedyny wektor \(x_0\),
 dla którego \(y_k=0\)
 dla wszystkich \(k\ge0\),
 to wektor zerowy.
\end{enumerate}
\end{twierdzenie}
\begin{proof}
Z równania stanu wynika
\(x_k=A^kx_0\),
a więc
\(y_k=CA^kx_0\).
Po ułożeniu pierwszych \(n\)
wektorów obserwacji
w jedną kolumnę dostajemy
\[
 \begin{bmatrix}
 y_0\\y_1\\\vdots\\y_{n-1}
 \end{bmatrix}
 =\mathcal O(C,A)x_0.
\]
Odwzorowanie
\(x_0\mapsto\mathcal O(C,A)x_0\)
jest iniektywne dokładnie wtedy,
gdy jego jądro jest zerowe.
To dowodzi równoważności
(i), (ii) i (iii).
Jeśli (i) zachodzi
i wszystkie \(y_k\) są zerowe,
to w szczególności
pierwszych \(n\) jest zerowych;
stąd \(x_0=0\) i otrzymujemy (iv).
Odwrotnie, jeśli
\(x_0\in\mathcal N(C,A)\),
to \(CA^kx_0=0\)
dla \(0\le k<n\).
Z twierdzenia Cayleya--Hamiltona
wynika ta sama równość
dla każdego \(k\ge n\),
gdyż \(A^k\) jest kombinacją
potęg \(I,A,\ldots,A^{n-1}\).
Zatem wszystkie obserwacje
tego stanu są zerowe.
Warunek (iv) wymusza
\(x_0=0\), a więc (i).
\end{proof}

\begin{lemat}[Niezmienniczość przestrzeni nieobserwowalnej]
\label{lem:lin25-unobservable-invariant}
Dla każdej pary \((C,A)\)
podprzestrzeń \(\mathcal N(C,A)\)
jest \(A\)-niezmiennicza
i jest największą
\(A\)-niezmienniczą podprzestrzenią
zawartą w \(\ker C\).
\end{lemat}
\begin{proof}
Jeżeli \(x\in\mathcal N(C,A)\),
to \(CA^kx=0\)
dla \(k=0,\ldots,n-1\).
Dla \(Ax\)
mamy \(CA^k(Ax)=CA^{k+1}x=0\)
dla \(0\le k<n-1\).
Także dla \(k=n-1\)
wyrażenie jest zerowe
na mocy twierdzenia
Cayleya--Hamiltona,
bo \(CA^nx\) jest kombinacją
\(CA^jx\) dla \(j<n\).
Zatem \(Ax\in\mathcal N(C,A)\).

Z definicji
\(\mathcal N(C,A)\subseteq\ker C\).
Jeśli \(L\subseteq\ker C\)
jest dowolną \(A\)-niezmienniczą
podprzestrzenią,
to dla każdego \(x\in L\)
wszystkie \(A^kx\)
należą do \(L\subseteq\ker C\).
Stąd \(CA^kx=0\)
dla wszystkich \(k\ge0\),
czyli \(x\in\mathcal N(C,A)\).
Otrzymujemy \(L\subseteq\mathcal N(C,A)\).
\end{proof}

\begin{twierdzenie}[Związek dualny z kryterium Kalmana]
\label{thm:lin25-duality}
Dla \(A\in\operatorname{Mat}_{n\times n}(\mathbb C)\)
i \(C\in\operatorname{Mat}_{p\times n}(\mathbb C)\)
para \((C,A)\) jest obserwowalna
dokładnie wtedy, gdy
para \((A^*,C^*)\)
jest sterowalna.
\end{twierdzenie}
\begin{proof}
Macierz sterowalności
pary \((A^*,C^*)\) ma postać
\[
 \begin{aligned}
 \mathcal K(A^*,C^*)
 &=
 \begin{bmatrix}
 C^*&A^*C^*&\cdots&(A^*)^{n-1}C^*
 \end{bmatrix}\\
 &=\mathcal O(C,A)^*,
 \end{aligned}
\]
ponieważ
\((CA^k)^*=(A^*)^kC^*\).
Macierz i jej sprzężenie
hermitowskie mają tę samą rangę.
Dlatego
\(\operatorname{rank}\mathcal K(A^*,C^*)=n\)
jest równoważne
\(\operatorname{rank}\mathcal O(C,A)=n\).
\end{proof}

\begin{twierdzenie}[Obserwowalność w czasie ciągłym]
\label{thm:lin25-cont-observe}
Dla układu
\(x'(t)=Ax(t)\),
\(y(t)=Cx(t)\),
gdzie \(x(0)=x_0\),
następujące warunki
są równoważne:
\begin{enumerate}[label=\textup{(\roman*)}]
 \item \(\operatorname{rank}\mathcal O(C,A)=n\);
 \item dla każdego \(T>0\)
 odwzorowanie
 \(x_0\mapsto(y(t))_{0\le t\le T}\)
 jest iniektywne;
 \item istnieje \(T>0\),
 dla którego to odwzorowanie
 jest iniektywne.
\end{enumerate}
\end{twierdzenie}
\begin{proof}
Mamy \(y(t)=Ce^{tA}x_0\).
Jeśli obserwacje dwóch stanów
początkowych są jednakowe
na \([0,T]\),
to dla ich różnicy \(v\)
zachodzi \(Ce^{tA}v=0\).
Różniczkując w \(t=0\)
dla \(k=0,\ldots,n-1\),
otrzymujemy
\(CA^kv=0\), czyli
\(\mathcal O(C,A)v=0\).
Jeśli (i) zachodzi,
to \(v=0\);
dowodzi to (ii).
Implikacja (ii) \(\Rightarrow\) (iii)
jest oczywista.

Załóżmy (iii).
Jeżeli
\(v\in\mathcal N(C,A)\),
to z twierdzenia
Cayleya--Hamiltona
\(CA^kv=0\)
dla wszystkich \(k\ge0\).
Stąd przez szereg wykładniczy
\(Ce^{tA}v=0\)
dla każdego \(t\ge0\).
Iniektywność dla pewnego \(T\)
wymusza \(v=0\).
Zatem \(\mathcal N(C,A)=\{0\}\),
co jest (i).
\end{proof}

\begin{przyklad}[Sterowalność i obserwowalność to różne własności]
\label{prz:lin25-dual-examples}
Niech
\[
 A=\begin{pmatrix}0&1\\0&0\end{pmatrix}.
\]
Dla \(B=e_2\)
para \((A,B)\) jest sterowalna,
bo \(\mathcal K(A,e_2)\)
ma kolumny \(e_2,e_1\).
Jednocześnie dla \(C=e_2^{\mathsf T}\)
mamy
\(CA=(0,0)\),
więc
\[
 \mathcal O(e_2^{\mathsf T},A)
 =\begin{pmatrix}0&1\\0&0\end{pmatrix}
\]
ma rangę \(1\):
druga para nie jest obserwowalna.

Z drugiej strony
dla \(B=e_1\)
mamy \(AB=0\),
czyli rangę sterowalności \(1\).
Dla \(C=e_1^{\mathsf T}\)
otrzymujemy
\[
 \mathcal O(e_1^{\mathsf T},A)
 =\begin{pmatrix}1&0\\0&1\end{pmatrix},
\]
więc para \((e_1^{\mathsf T},A)\)
jest obserwowalna.
\end{przyklad}

\section{Spektralne kryterium PBH}
\label{sec:lin25-pbh}

\begin{twierdzenie}[Kryterium Popowa--Belevitcha--Hautusa]
\label{thm:lin25-pbh}
Niech \(n,m\ge1\) będą liczbami całkowitymi,
\(A\in\operatorname{Mat}_{n\times n}(\mathbb C)\)
oraz \(B\in\operatorname{Mat}_{n\times m}(\mathbb C)\).
Para \((A,B)\)
jest sterowalna dokładnie wtedy,
gdy dla każdego
\(\lambda\in\mathbb C\)
zachodzi
\[
 \operatorname{rank}
 \begin{bmatrix}\lambda I_n-A&B\end{bmatrix}=n.
\]
Wystarczy sprawdzać wartości
\(\lambda\in\sigma(A)\).
Równoważnie nie istnieją
\(\lambda\in\mathbb C\)
i niezerowy wiersz
\(\varphi\in(\mathbb C^n)^*\)
takie, że
\[
 \varphi A=\lambda\varphi,
 \qquad \varphi B=0.
\]
\end{twierdzenie}
\begin{proof}
Macierz o \(n\) wierszach
ma rangę mniejszą od \(n\)
dokładnie wtedy, gdy istnieje
niezerowy wiersz \(\varphi\)
zerujący iloczyn z nią.
Dla macierzy
\([\lambda I_n-A\ \ B]\)
warunki te mają postać
\(\varphi A=\lambda\varphi\)
i \(\varphi B=0\).
Pierwszy z nich oznacza,
że \(\lambda\) jest wartością
własną \(A\), bo macierze
\(A\) i \(A^{\mathsf T}\)
mają ten sam wielomian
charakterystyczny.
Dlatego wystarczy
sprawdzać \(\sigma(A)\).

Jeśli istnieje taka para
\((\lambda,\varphi)\),
to dla każdego \(k\ge0\)
\[
 \varphi A^kB
 =\lambda^k\varphi B=0.
\]
Wobec tego
\(\varphi\mathcal K(A,B)=0\),
a ponieważ \(\varphi\ne0\),
ranga \(\mathcal K(A,B)\)
jest mniejsza od \(n\).
Para nie jest sterowalna.

Odwrotnie, jeśli
\(\mathcal R(A,B)\ne\mathbb C^n\),
to jej anihilator
\[
 \mathcal R(A,B)^\circ
 =\{\varphi\in(\mathbb C^n)^*:
        \varphi r=0\text{ dla }r\in\mathcal R(A,B)\}
\]
jest niezerową przestrzenią
skończenie wymiarową.
Z Lematu~\ref{lem:lin25-invariant}
wynika, że \(\mathcal R(A,B)\)
jest \(A\)-niezmiennicza.
Dlatego
\(\varphi\mapsto\varphi A\)
zachowuje anihilator:
dla \(r\in\mathcal R(A,B)\)
mamy
\((\varphi A)r=\varphi(Ar)=0\).
Każdy endomorfizm niezerowej
przestrzeni zespolonej
ma wektor własny.
Istnieją zatem
\(\varphi\ne0\) w anihilatorze
i \(\lambda\in\mathbb C\)
takie, że
\(\varphi A=\lambda\varphi\).
Ponieważ \(\operatorname{im}B
 \subseteq\mathcal R(A,B)\),
mamy też \(\varphi B=0\).
Otrzymaliśmy awarię
warunku rangowego PBH,
co dowodzi twierdzenia.
\end{proof}

\begin{wniosek}[PBH dla obserwowalności]
\label{cor:lin25-pbh-observe}
Dla
\(A\in\operatorname{Mat}_{n\times n}(\mathbb C)\)
i
\(C\in\operatorname{Mat}_{p\times n}(\mathbb C)\)
para \((C,A)\)
jest obserwowalna dokładnie wtedy,
gdy dla każdego
\(\lambda\in\sigma(A)\)
zachodzi
\[
 \operatorname{rank}
 \begin{bmatrix}
 \lambda I_n-A\\C
 \end{bmatrix}=n.
\]
Równoważnie nie istnieje
niezerowy \(v\in\mathbb C^n\)
spełniający
\(Av=\lambda v\)
i \(Cv=0\).
\end{wniosek}
\begin{proof}
Macierz z tezy ma jądro
\(\ker(\lambda I_n-A)\cap\ker C\).
Jej ranga wynosi \(n\)
dokładnie wtedy, gdy
przecięcie jest zerowe.
Załóżmy, że istnieje
niezerowy \(v\)
z \(Av=\lambda v\) i \(Cv=0\).
Wtedy \(CA^kv=\lambda^kCv=0\)
dla każdego \(k\),
więc \(v\in\mathcal N(C,A)\);
para nie jest obserwowalna.

Odwrotnie, jeśli
\(\mathcal N(C,A)\ne\{0\}\),
to z
Lematu~\ref{lem:lin25-unobservable-invariant}
jest ona niezerową
\(A\)-niezmienniczą
podprzestrzenią.
Ograniczenie \(A|_{\mathcal N}\)
ma wektor własny \(v\ne0\).
Ponieważ
\(\mathcal N(C,A)\subseteq\ker C\),
otrzymujemy \(Cv=0\),
a zatem brak wymaganego
warunku rangi.
\end{proof}

\begin{przyklad}[Zawodność testu opartego tylko na wyznaczniku \(A\)]
\label{prz:lin25-pbh}
Niech
\[
 A=\operatorname{diag}(1,2),
 \qquad B=\begin{pmatrix}1\\0\end{pmatrix}.
\]
Macierz \(A\) jest odwracalna,
ale
\[
 \mathcal K(A,B)
 =\begin{pmatrix}1&1\\0&0\end{pmatrix}
\]
ma rangę \(1\).
Dla \(\lambda=2\)
wiersz \(\varphi=e_2^{\mathsf T}\)
spełnia
\(\varphi A=2\varphi\)
i \(\varphi B=0\).
Zatem kryterium PBH nie jest
spełnione. Odwracalność \(A\)
nie implikuje sterowalności,
bo sterowanie nie działa
na drugi kierunek własny.
\end{przyklad}

\section{Rozkład Kalmana}
\label{sec:lin25-kalman-decomposition}

\begin{twierdzenie}[Rozkład na część osiągalną i nieosiągalną]
\label{thm:lin25-reachable-split}
Niech
\(\mathcal R=\mathcal R(A,B)\)
i \(r=\dim_{\mathbb C}\mathcal R\).
Istnieje baza przestrzeni
\(\mathbb C^n\), której pierwsze
\(r\) wektorów tworzy bazę
\(\mathcal R\), i w której
macierze układu mają postać
\[
 A=\begin{pmatrix}
 A_{\mathrm r}&D\\0&A_{\mathrm u}
 \end{pmatrix},
 \qquad
 B=\begin{pmatrix}B_{\mathrm r}\\0\end{pmatrix}.
\]
Występujące bloki mają wymiary
\(r\times r\), \(r\times(n-r)\),
\((n-r)\times(n-r)\)
oraz \(r\times m\).
Para \((A_{\mathrm r},B_{\mathrm r})\)
jest sterowalna
w przestrzeni \(\mathbb C^r\).
Operator \(A_{\mathrm u}\)
jest macierzą endomorfizmu
indukowanego przez \(A\)
na ilorazie
\(\mathbb C^n/\mathcal R\);
jego podobieństwa, a więc
i wartości własne, nie zależą
od wyboru dopełnienia.

Jeżeli w układzie
\(x'=Ax+Bu\) albo
\(x_{k+1}=Ax_k+Bu_k\)
zastosujemy liniowe
sprzężenie \(u=-Kx\),
to indukowany operator
zamkniętej pętli na ilorazie
\(\mathbb C^n/\mathcal R\)
nadal jest równy
operatorowi indukowanemu
przez \(A\).
\end{twierdzenie}
\begin{proof}
Z
Lematu~\ref{lem:lin25-invariant}
\(\mathcal R\) jest
\(A\)-niezmiennicza
i zawiera \(\operatorname{im}B\).
Wybierzmy bazę \(\mathcal R\)
i rozszerzmy ją do bazy
całej przestrzeni.
Niezmienniczość
oznacza, że obrazy pierwszych
\(r\) wektorów bazowych
nie mają współrzędnych
w ostatnich \(n-r\) kierunkach,
co daje lewy dolny blok
zerowy w macierzy \(A\).
Warunek
\(\operatorname{im}B\subseteq\mathcal R\)
daje dolny blok zerowy macierzy \(B\).

Dla każdego \(k\ge0\)
wektory \(A^kBu\) należą
do \(\mathcal R\) i mają
w wybranej bazie współrzędne
\(((A_{\mathrm r})^kB_{\mathrm r}u,0)\).
Ponieważ ich suma liniowa
generuje \(\mathcal R\),
pierwsze \(r\) współrzędnych
generują \(\mathbb C^r\).
Stąd
\(\operatorname{rank}
 \mathcal K(A_{\mathrm r},B_{\mathrm r})=r\)
na mocy
Lematu~\ref{lem:lin25-invariant}
zastosowanego w wymiarze \(r\).

Na ilorazie
\(\mathbb C^n/\mathcal R\)
odwzorowanie
\([x]\mapsto[Ax]\)
jest dobrze określone,
bo \(\mathcal R\)
jest \(A\)-niezmiennicza.
W wybranej bazie
jego macierzą jest
\(A_{\mathrm u}\).
Inne dopełnienie zmienia
jedynie bazę ilorazu,
stąd podobieństwo.

Jeśli \(u=-Kx\),
nową macierzą stanu
jest \(A-BK\).
Ponieważ \(\operatorname{im}BK
 \subseteq\operatorname{im}B
 \subseteq\mathcal R\),
dla dowolnego \(x\)
mamy
\([(A-BK)x]=[Ax]\)
w ilorazie.
Wynik jest jednakowy
dla czasu ciągłego
i dyskretnego.
Przypadki \(r=0\) i \(r=n\)
rozumiemy przez pominięcie
pustych bloków.
\end{proof}

\begin{twierdzenie}[Rozkład przestrzeni według sterowalności i obserwowalności]
\label{thm:lin25-kalman-decomposition}
Niech
\(\mathcal R=\mathcal R(A,B)\)
i \(\mathcal N=\mathcal N(C,A)\).
Istnieją podprzestrzenie
\(U,R_0,N_0,W\subseteq\mathbb C^n\)
takie, że
\[
 \begin{gathered}
 U=\mathcal R\cap\mathcal N,\qquad
 \mathcal R=U\oplus R_0,\qquad
 \mathcal N=U\oplus N_0,\\
 \mathbb C^n
 =U\oplus R_0\oplus N_0\oplus W,
 \end{gathered}
\]
oraz w bazie dostosowanej
do tego rozkładu
macierze mają postać
\[
 A=\begin{pmatrix}
 A_U&*&*&*\\
 0&A_{R_0}&0&*\\
 0&0&A_{N_0}&*\\
 0&0&0&A_W
 \end{pmatrix},
 \quad
 B=\begin{pmatrix}B_U\\B_{R_0}\\0\\0\end{pmatrix},
 \quad
 C=\begin{pmatrix}0&C_{R_0}&0&C_W\end{pmatrix}.
\]
Każda gwiazdka oznacza
blok odpowiedniego rozmiaru,
który nie musi być zerowy.
Składniki odpowiadają
kolejno części osiągalnej
nieobserwowalnej,
osiągalnej obserwowalnej,
nieosiągalnej nieobserwowalnej
i nieosiągalnej obserwowalnej,
z interpretacją odpowiednich
podilorazów.
\end{twierdzenie}
\begin{proof}
Z Lematu~\ref{lem:lin25-invariant}
\(\mathcal R\) jest
\(A\)-niezmiennicza.
Z
Lematu~\ref{lem:lin25-unobservable-invariant}
to samo zachodzi dla
\(\mathcal N\).
Stąd ich przecięcie
\(U\) i suma
\(\mathcal R+\mathcal N\)
również są \(A\)-niezmiennicze.

Wybierzmy bazę \(U\).
Rozszerzmy ją osobno
do bazy \(\mathcal R\)
oraz do bazy \(\mathcal N\).
Niech \(R_0\) i \(N_0\)
będą przestrzeniami
generowanymi przez dodane
wektory. Z faktu
\(\mathcal R\cap\mathcal N=U\)
wynika, że
\(U\oplus R_0\oplus N_0
 =\mathcal R+\mathcal N\)
jest sumą prostą.
Uzupełnijmy jej bazę
do bazy całej przestrzeni,
otrzymując \(W\).

Rozpatrzmy kolumny
macierzy \(A\)
w kolejności
\(U,R_0,N_0,W\).
Ponieważ \(A(U)\subseteq U\),
pierwszy blok kolumn
ma niezerowe współczynniki
wyłącznie w pierwszym
bloku wierszy.
Z \(A(\mathcal R)\subseteq\mathcal R\)
wynika, że kolumny
odpowiadające \(R_0\)
mogą mieć niezerowe składowe
wyłącznie w \(U\) i \(R_0\).
Podobnie kolumny \(N_0\)
mogą mieć niezerowe składowe
wyłącznie w \(U\) i \(N_0\).
Dla \(W\) nie ma ograniczenia.
Daje to wskazaną postać \(A\).

Obraz \(B\) leży w \(\mathcal R\),
więc jego współrzędne
w \(N_0,W\) są zerowe.
Natomiast \(C\) znika
na całej \(\mathcal N\),
a więc na \(U\) i \(N_0\);
pozostałe bloki nie muszą
się zerować. Otrzymujemy
podane postacie \(B,C\).

Nazwy czterech części
należy rozumieć przez
ich rolę w filtracji
\[
 0\subseteq U
 \subseteq\mathcal R,\mathcal N
 \subseteq\mathcal R+\mathcal N
 \subseteq\mathbb C^n;
\]
wybrane dopełnienia
nie muszą same być
\(A\)-niezmiennicze.
Na przykład obserwowalna
część osiągalna jest
reprezentowana przez
\(\mathcal R/U\),
nie przez wyróżnioną
kanonicznie podprzestrzeń \(R_0\).
\end{proof}

\begin{przyklad}[Rozkład z niesterowalnym, stabilnym kierunkiem]
\label{prz:lin25-reachable-split}
Niech
\[
 A=\begin{pmatrix}1&0\\0&-2\end{pmatrix},
 \qquad B=\begin{pmatrix}1\\0\end{pmatrix},
 \qquad C=\begin{pmatrix}1&0\end{pmatrix}.
\]
Mamy
\(\mathcal R=\operatorname{span}(e_1)\)
oraz
\(\mathcal N=\operatorname{span}(e_2)\).
Rozkład Kalmana ma zatem
dwa niezerowe składniki:
\(R_0=\operatorname{span}(e_1)\)
i \(N_0=\operatorname{span}(e_2)\).
Sterowanie działa wyłącznie
na pierwszą współrzędną,
a obserwacja nie widzi drugiej.
Stan w drugim kierunku
zanika jednak samoczynnie
w czasie ciągłym, gdyż
\(x_2'(t)=-2x_2(t)\).
Później zobaczymy, że
ta para jest stabilizowalna
i wykrywalna, mimo braku
pełnej sterowalności
i obserwowalności.
\end{przyklad}

\section{Rozmieszczanie stabilnych wartości własnych}
\label{sec:lin25-feedback}

\begin{lemat}[Konstrukcja stabilizującego sprzężenia dla pary sterowalnej]
\label{lem:lin25-feedback-placement}
Niech
\(H\in\operatorname{Mat}_{r\times r}(\mathbb C)\)
i \(D\in\operatorname{Mat}_{r\times m}(\mathbb C)\),
gdzie \(r,m\ge1\),
a para \((H,D)\)
jest sterowalna.
Niech \(\Omega\subseteq\mathbb C\)
będzie niepustym zbiorem
otwartym.
Wówczas istnieją
\emph{różne} punkty
\(\mu_1,\ldots,\mu_r
 \in\Omega\setminus\sigma(H)\)
i macierz
\(F\in\operatorname{Mat}_{m\times r}(\mathbb C)\)
takie, że wartościami własnymi
macierzy \(H-DF\)
są dokładnie
\(\mu_1,\ldots,\mu_r\).
\end{lemat}
\begin{proof}
Dla
\(\mu\notin\sigma(H)\)
oznaczmy
\[
 E_\mu=\operatorname{im}(H-\mu I_r)^{-1}D
 \subseteq\mathbb C^r.
\]
Wykażemy, że dla
\emph{każdej} właściwej
podprzestrzeni
\(W\subsetneq\mathbb C^r\)
oraz dowolnego skończonego
zbioru \(S\subseteq\mathbb C\)
można wybrać
\(\mu\in\Omega\setminus
 (\sigma(H)\cup S)\)
takie, że
\(E_\mu\not\subseteq W\).

Gdyby tak nie było,
to wszystkie \(E_\mu\)
dla \(\mu\) z niepustego
otwartego zbioru
\(\Omega\setminus(\sigma(H)\cup S)\)
leżałyby w \(W\).
Wybierzmy niezerowy
funkcjonał liniowy
\(\varphi\in(\mathbb C^r)^*\)
zerujący \(W\);
istnieje on, bo \(W\)
jest właściwa.
Dla każdego takiego \(\mu\)
otrzymalibyśmy
\[
 \varphi(H-\mu I_r)^{-1}D=0.
\]
Każda współrzędna tego
wiersza jest funkcją wymierną
zmiennej \(\mu\), gdyż
\[
 (H-\mu I_r)^{-1}
 =\frac{\operatorname{adj}(H-\mu I_r)}
         {\det(H-\mu I_r)}.
\]
Funkcja wymierna znikająca
na niepustym otwartym zbiorze
jest tożsamościowo zerowa:
jej licznik jest wielomianem
o nieskończenie wielu zerach.
Dla dostatecznie dużych
\(|\mu|\) możemy rozwinąć
odwrotność w szereg geometryczny:
\[
 (H-\mu I_r)^{-1}
 =-\sum_{k=0}^{\infty}
      \mu^{-k-1}H^k.
\]
Z równości
\(\varphi(H-\mu I_r)^{-1}D=0\)
wynika przez jednoznaczność
współczynników zbieżnego
szeregu potęgowego
w \(1/\mu\),
że \(\varphi H^kD=0\)
dla wszystkich \(k\ge0\).
W szczególności
\(\varphi\mathcal K(H,D)=0\),
co przeczy sterowalności
pary \((H,D)\).
Zatem wskazany wybór
\(\mu\) istnieje.

Zaczynamy od \(W_0=\{0\}\).
Stosując powyższą własność
indukcyjnie, wybieramy
różne
\(\mu_1,\ldots,\mu_r\)
z \(\Omega\setminus\sigma(H)\)
oraz wektory
\(u_j\in\mathbb C^m\)
takie, że
\[
 v_j=(H-\mu_j I_r)^{-1}Du_j
 \notin\operatorname{span}
       (v_1,\ldots,v_{j-1})
 \quad(1\le j\le r).
\]
Dla \(j=1\) prawa strona
oznacza \(\{0\}\).
Wektory \(v_1,\ldots,v_r\)
są liniowo niezależne,
a więc tworzą bazę
\(\mathbb C^r\).
Określmy odwzorowanie
liniowe \(F\)
jednoznacznie przez
\(Fv_j=u_j\).
Z definicji \(v_j\)
wynika
\[
 Hv_j-\mu_jv_j=Du_j,
 \qquad
 (H-DF)v_j=\mu_jv_j.
\]
Zatem w bazie
\(v_1,\ldots,v_r\)
macierz \(H-DF\)
jest diagonalna
z elementami
\(\mu_1,\ldots,\mu_r\).
To kończy dowód.
\end{proof}

\begin{uwaga}[Co zostało rzeczywiście wykazane]
\label{rem:lin25-poles}
Lemat~\ref{lem:lin25-feedback-placement}
pozwala \emph{wybrać}
różne wartości własne
sprzężenia w dowolnym
otwartym obszarze \(\Omega\).
Nie stwierdza,
że dowolny z góry zadany
ciąg wartości własnych
można uzyskać
przez to samo sprzężenie:
pełne twierdzenie
o rozmieszczaniu biegunów
jest silniejszym wynikiem.
Do udowodnienia
stabilizowalności
wystarczy podana konstrukcja.
\end{uwaga}

\section{Stabilizowalność}
\label{sec:lin25-stabilizability}

\begin{definicja}[Stabilizowalność w obu rodzajach czasu]
\label{def:lin25-stabilizability}
Parę \((A,B)\)
nazywamy \emph{stabilizowalną
w czasie ciągłym},
jeżeli istnieje
\(K\in\operatorname{Mat}_{m\times n}(\mathbb C)\)
takie, że każda wartość własna
\(A-BK\) ma ujemną część
rzeczywistą.
Nazywamy ją
\emph{stabilizowalną
w czasie dyskretnym},
jeżeli istnieje taka macierz \(K\),
dla której
\(\rho(A-BK)<1\).
W obu definicjach \(K\)
jest stałym sprzężeniem
zwrotnym od stanu:
\(u=-Kx\).
\end{definicja}

\begin{twierdzenie}[Kryterium stabilizowalności]
\label{thm:lin25-stabilizability}
Dla dowolnej pary
\(A\in\operatorname{Mat}_{n\times n}(\mathbb C)\),
\(B\in\operatorname{Mat}_{n\times m}(\mathbb C)\)
następujące warunki
są równoważne:
\begin{enumerate}[label=\textup{(\roman*)}]
 \item \((A,B)\)
 jest stabilizowalna
 w czasie ciągłym;
 \item dla każdej wartości własnej
 \(\lambda\in\sigma(A)\)
 z \(\operatorname{Re}\lambda\ge0\)
 zachodzi
 \[
  \operatorname{rank}
  \begin{bmatrix}\lambda I_n-A&B\end{bmatrix}=n;
 \]
 \item w rozkładzie
 z Twierdzenia~\ref{thm:lin25-reachable-split}
 wszystkie wartości własne
 \(A_{\mathrm u}\)
 mają ujemną część rzeczywistą.
\end{enumerate}
Analogicznie następujące
warunki są równoważne:
\begin{enumerate}[label=\textup{(\roman*)},
                  start=4]
 \item \((A,B)\)
 jest stabilizowalna
 w czasie dyskretnym;
 \item warunek rangi
 \[
  \operatorname{rank}
  \begin{bmatrix}\lambda I_n-A&B\end{bmatrix}=n
 \]
 zachodzi dla każdego
 \(\lambda\in\sigma(A)\)
 z \(|\lambda|\ge1\);
 \item w rozkładzie
 z Twierdzenia~\ref{thm:lin25-reachable-split}
 wszystkie wartości własne
 \(A_{\mathrm u}\)
 mają moduł mniejszy niż \(1\).
\end{enumerate}
Jeśli \((A,B)\) jest sterowalna,
to jest stabilizowalna
w obu rodzajach czasu.
\end{twierdzenie}
\begin{proof}
Ustalmy rozkład z
Twierdzenia~\ref{thm:lin25-reachable-split}:
\[
 A=\begin{pmatrix}A_{\mathrm r}&D\\0&A_{\mathrm u}\end{pmatrix},
 \qquad
 B=\begin{pmatrix}B_{\mathrm r}\\0\end{pmatrix}.
\]
Wykazaliśmy tam, że
\((A_{\mathrm r},B_{\mathrm r})\)
jest sterowalna.
Pokażemy najpierw,
że dla każdej
\(\lambda\in\mathbb C\)
\[
 \operatorname{rank}
 \begin{bmatrix}\lambda I_n-A&B\end{bmatrix}<n
 \quad\Longleftrightarrow\quad
 \lambda\in\sigma(A_{\mathrm u}).
\]
Jeśli ranga po lewej
jest mniejsza od \(n\),
to na mocy
Twierdzenia~\ref{thm:lin25-pbh}
istnieje niezerowy wiersz
\(\varphi=(\varphi_{\mathrm r},
\varphi_{\mathrm u})\)
taki, że
\(\varphi A=\lambda\varphi\)
i \(\varphi B=0\).
Porównując pierwsze
\(r\) kolumn, otrzymujemy
\[
 \varphi_{\mathrm r}A_{\mathrm r}
  =\lambda\varphi_{\mathrm r},
 \qquad
 \varphi_{\mathrm r}B_{\mathrm r}=0.
\]
Sterowalność części osiągalnej
i kryterium PBH wymuszają
\(\varphi_{\mathrm r}=0\).
Ponieważ \(\varphi\ne0\),
otrzymujemy
\(\varphi_{\mathrm u}\ne0\).
Porównanie pozostałych kolumn
daje
\(\varphi_{\mathrm u}A_{\mathrm u}
 =\lambda\varphi_{\mathrm u}\),
czyli \(\lambda\in\sigma(A_{\mathrm u})\).

Odwrotnie, jeśli
\(\lambda\) jest wartością własną
\(A_{\mathrm u}\),
to istnieje niezerowy
lewy wektor własny
\(\psi A_{\mathrm u}
 =\lambda\psi\).
Wiersz \(\varphi=(0,\psi)\)
spełnia
\(\varphi A=\lambda\varphi\)
oraz \(\varphi B=0\);
a więc ranga jest
mniejsza od \(n\).
Dowiedliśmy równoważności.
W przypadkach \(r=0\) i \(r=n\)
rozumowanie dotyczy
odpowiednio pełnego operatora
lub pustego widma ilorazu.

Jeżeli istnieje sprzężenie
\(K\), dla którego \(A-BK\)
jest stabilna w odpowiednim
rodzaju czasu, to
z
Twierdzenia~\ref{thm:lin25-reachable-split}
operator indukowany
przez \(A-BK\) na ilorazie
\(\mathbb C^n/\mathcal R\)
jest równy operatorowi
indukowanemu przez \(A\).
W odpowiedniej bazie
macierz zamkniętej pętli
ma postać górnotrójkątną
z diagonalnym blokiem
\(A_{\mathrm u}\).
Jej wielomian
charakterystyczny jest
iloczynem wielomianów
bloków diagonalnych,
więc każde
\(\lambda\in\sigma(A_{\mathrm u})\)
musi należeć do
obszaru stabilności.
Dowodzi to (i) \(\Rightarrow\) (iii)
oraz (iv) \(\Rightarrow\) (vi).

Załóżmy (iii).
Jeśli \(r>0\),
zastosujmy
Lemat~\ref{lem:lin25-feedback-placement}
do sterowalnej pary
\((A_{\mathrm r},B_{\mathrm r})\)
oraz otwartego obszaru
\(\Omega_{\mathrm c}
 =\{\mu\in\mathbb C:
           \operatorname{Re}\mu<0\}\).
Otrzymujemy macierz \(F\)
taką, że wszystkie wartości
własne \(A_{\mathrm r}-B_{\mathrm r}F\)
mają ujemne części rzeczywiste.
Jeśli \(r=0\), ten krok pomijamy.
W bazie rozkładu Kalmana
wybierzmy sprzężenie
\(K=(F\;\;0)\).
Wówczas
\[
 A-BK=
 \begin{pmatrix}
 A_{\mathrm r}-B_{\mathrm r}F&D\\
 0&A_{\mathrm u}
 \end{pmatrix}.
\]
Każda wartość własna
obu bloków diagonalnych
ma ujemną część rzeczywistą,
więc \(A-BK\) jest stabilna
w czasie ciągłym.
Dowodzi to (iii) \(\Rightarrow\) (i).

Analogicznie,
jeśli zachodzi (vi),
zastosujmy ten sam lemat
z obszarem
\(\Omega_{\mathrm d}
 =\{\mu\in\mathbb C:|\mu|<1\}\).
Otrzymamy
\(\rho(A_{\mathrm r}-B_{\mathrm r}F)<1\)
i dzięki górnotrójkątnej
postaci zamkniętej pętli
\(\rho(A-BK)<1\).
Zatem (vi) \(\Rightarrow\) (iv).

Wykazana na początku dowodu
równoważność awarii PBH
i \(\lambda\in\sigma(A_{\mathrm u})\)
daje jednocześnie
(ii) \(\Leftrightarrow\) (iii)
i (v) \(\Leftrightarrow\) (vi).
Jeśli para jest sterowalna,
to \(\mathcal R=\mathbb C^n\);
operator na zerowym ilorazie
nie ma wartości własnych,
więc (iii) i (vi)
są spełnione automatycznie.
\end{proof}

\begin{przyklad}[Stabilizowalność bez sterowalności]
\label{prz:lin25-stabilizable}
Niech
\[
 A=\begin{pmatrix}1&0\\0&-2\end{pmatrix},
 \qquad B=\begin{pmatrix}1\\0\end{pmatrix}.
\]
W Przykładzie~\ref{prz:lin25-reachable-split}
wykazaliśmy, że
\(\mathcal R=\operatorname{span}(e_1)\).
Macierz na nieosiągalnym
ilorazie jest skalarna
\(A_{\mathrm u}=(-2)\),
więc para jest stabilizowalna
w czasie ciągłym.
Rzeczywiście wybór
\(K=(3\;\;0)\)
daje
\[
 A-BK=\operatorname{diag}(-2,-2),
\]
której obie wartości własne
mają ujemne części rzeczywiste.

Nie jest to jednak
stabilizowalność w czasie
dyskretnym przy tym samym \(A\):
nieosiągalna wartość własna
\(-2\) ma moduł \(2>1\).
Żadna macierz \(K\)
nie może przenieść tej
wartości własnej do dysku
jednostkowego.
Różnica między stabilnością
ciągłą a dyskretną
jest tu istotna.
\end{przyklad}

\begin{przyklad}[Jednoczesna porażka obu kryteriów]
\label{prz:lin25-not-stabilizable}
Dla
\[
 A=\operatorname{diag}(1,2),
 \qquad B=e_1
\]
nieosiągalny kierunek \(e_2\)
ma wartość własną \(2\).
Dla dowolnego
\(K=(k_1\;\;k_2)\)
macierz
\[
 A-BK=
 \begin{pmatrix}1-k_1&-k_2\\0&2\end{pmatrix}
\]
ma nadal wartość własną \(2\).
Zatem para nie jest
stabilizowalna ani w czasie
ciągłym, ani w dyskretnym.
Nie pomaga nawet pełna
swoboda doboru obu
współczynników \(K\):
sterowanie nie zmienia
dynamiki ilorazu.
\end{przyklad}

\section{Wykrywalność}
\label{sec:lin25-detectability}

\begin{definicja}[Wykrywalność]
\label{def:lin25-detectability}
Dla pary \((C,A)\)
mówimy o \emph{wykrywalności
w czasie ciągłym},
jeśli istnieje
\(L\in\operatorname{Mat}_{n\times p}(\mathbb C)\)
takie, że
\(A-LC\) ma wszystkie
wartości własne
o ujemnych częściach rzeczywistych.
Para jest \emph{wykrywalna
w czasie dyskretnym},
jeśli istnieje takie \(L\),
dla którego \(\rho(A-LC)<1\).
Macierz \(L\) nazywamy
\emph{wzmocnieniem obserwatora}.
\end{definicja}

\begin{twierdzenie}[Spektralne kryteria wykrywalności]
\label{thm:lin25-detectability}
Para
\(C\in\operatorname{Mat}_{p\times n}(\mathbb C)\),
\(A\in\operatorname{Mat}_{n\times n}(\mathbb C)\)
jest wykrywalna w czasie ciągłym
dokładnie wtedy, gdy
dla każdej wartości własnej
\(\lambda\in\sigma(A)\)
z \(\operatorname{Re}\lambda\ge0\)
zachodzi
\[
 \operatorname{rank}
 \begin{bmatrix}\lambda I_n-A\\C\end{bmatrix}=n.
\]
Jest wykrywalna
w czasie dyskretnym
dokładnie wtedy, gdy
ten sam warunek zachodzi
dla wszystkich
\(\lambda\in\sigma(A)\)
z \(|\lambda|\ge1\).
Każda para obserwowalna
jest wykrywalna
w obu rodzajach czasu.
\end{twierdzenie}
\begin{proof}
Z
Twierdzenia~\ref{thm:lin25-duality}
para \((C,A)\)
jest obserwowalna
dokładnie wtedy, gdy
para \((A^*,C^*)\)
jest sterowalna.
Tę samą operację sprzężenia
zastosujemy do stabilizowalności.
Dla dowolnego
\(L\in\operatorname{Mat}_{n\times p}(\mathbb C)\)
mamy
\[
 (A-LC)^*=A^*-C^*L^*.
\]
Wartości własne
macierzy sprzężonej
są sprzężeniami zespolonymi
wartości własnych
macierzy wyjściowej.
Dlatego \(A-LC\)
jest stabilna w czasie
ciągłym lub dyskretnym
dokładnie wtedy, gdy
odpowiednio stabilna
jest \(A^*-C^*L^*\).
Wobec tego wykrywalność
\((C,A)\) jest równoważna
stabilizowalności
\((A^*,C^*)\).

Z
Twierdzenia~\ref{thm:lin25-stabilizability}
warunkiem tej ostatniej
jest pełna ranga
\[
 \begin{bmatrix}
 \overline{\lambda}I_n-A^*&C^*
 \end{bmatrix}
\]
dla każdej wartości własnej
\(\overline{\lambda}\in\sigma(A^*)\),
której część rzeczywista
jest nieujemna
lub której moduł
jest co najmniej \(1\).
Po sprzężeniu całej macierzy
otrzymujemy pionową
macierz
\(\left[\begin{smallmatrix}
 \lambda I_n-A\\C
 \end{smallmatrix}\right]\).
Macierz i jej sprzężenie
hermitowskie mają tę samą rangę,
więc otrzymujemy dokładnie
warunki z tezy.
Na koniec z
Wniosku~\ref{cor:lin25-pbh-observe}
para obserwowalna
spełnia warunek rangowy
dla wszystkich
wartości własnych,
w szczególności
dla niestabilnych.
\end{proof}

\begin{przyklad}[Niewidoczny, ale stabilny stan]
\label{prz:lin25-detectable}
Dla
\[
 A=\operatorname{diag}(1,-2),
 \qquad C=\begin{pmatrix}1&0\end{pmatrix}
\]
podprzestrzeń nieobserwowalna
jest równa
\(\operatorname{span}(e_2)\),
bo \(\mathcal O(C,A)\)
ma rangę \(1\).
Nieobserwowalna wartość
własna \(-2\)
jest stabilna w czasie ciągłym.
Zatem para jest wykrywalna
w czasie ciągłym.
Na przykład dla
\(L=(3,0)^{\mathsf T}\)
otrzymujemy
\(A-LC=\operatorname{diag}(-2,-2)\).

W czasie dyskretnym
wartość własna \(-2\)
jest niestabilna.
Żadne \(L\)
nie może jej przesunąć:
\[
 A-LC=
 \begin{pmatrix}1-l_1&0\\-l_2&-2\end{pmatrix}
\]
zawsze ma wartość własną \(-2\).
Wykrywalność nie zachodzi
w tym rodzaju czasu.
\end{przyklad}

\begin{stwierdzenie}[Dynamika błędu obserwatora]
\label{prop:lin25-observer}
Niech \(u(t)\)
będzie ustalonym
przedziałami ciągłym
sterowaniem, a stan
układu ciągłego
niech spełnia
\(x'=Ax+Bu\), \(y=Cx\).
Dla dowolnej macierzy
\(L\in\operatorname{Mat}_{n\times p}(\mathbb C)\)
rozważmy obserwator
\[
 \widehat x'
 =A\widehat x+Bu
  +L(y-C\widehat x).
\]
Błąd \(e=x-\widehat x\)
spełnia równanie
\(e'=(A-LC)e\).
Jeżeli \(A-LC\)
jest spektralnie stabilna
w czasie ciągłym,
to \(e(t)\to0\)
wykładniczo dla każdego
początkowego błędu.

W czasie dyskretnym,
przy \(x_{k+1}=Ax_k+Bu_k\),
\(y_k=Cx_k\) oraz
\[
 \widehat x_{k+1}
 =A\widehat x_k+Bu_k
  +L(y_k-C\widehat x_k),
\]
otrzymujemy
\(e_{k+1}=(A-LC)e_k\);
jeżeli \(\rho(A-LC)<1\),
to \(e_k\to0\) geometrycznie.
\end{stwierdzenie}
\begin{proof}
Odejmując równanie obserwatora
od równania stanu,
otrzymujemy
\[
 \begin{aligned}
 e'
 &=Ax+Bu-
    (A\widehat x+Bu+LCx-LC\widehat x)\\
 &=(A-LC)(x-\widehat x)
 =(A-LC)e.
 \end{aligned}
\]
Z
Rozdziału~\ref{ch:lin22-sylvester}
wynika, że jeśli każda wartość
własna \(A-LC\) ma ujemną
część rzeczywistą, to macierz
\(e^{t(A-LC)}\)
zanika wykładniczo;
a więc
\(e(t)=e^{t(A-LC)}e(0)\to0\).
W czasie dyskretnym
tę samą operację odejmowania
wykonujemy bez pochodnych:
\[
 e_{k+1}
 =Ax_k+Bu_k-
 (A\widehat x_k+Bu_k+LCx_k-LC\widehat x_k)
 =(A-LC)e_k.
\]
Z
Lematu~\ref{lem:lin23-powers}
otrzymujemy
\(e_k=(A-LC)^ke_0\to0\)
i oszacowanie geometryczne.
\end{proof}

\section{Energia sterowania i gramian dyskretny}
\label{sec:lin25-discrete-gramian}

\begin{stwierdzenie}[Gramian na skończonym horyzoncie]
\label{prop:lin25-discrete-gramian}
Niech \(N\ge1\) będzie
liczbą całkowitą.
Dla pary \((A,B)\)
określmy
\[
 W_N^{\mathrm d}=
 \sum_{j=0}^{N-1}A^jBB^*(A^*)^j.
\]
Jest to macierz hermitowska
i nieujemnie określona.
Dla każdego \(N\)
zachodzi
\[
 W_N^{\mathrm d}=
 \mathcal K_N\mathcal K_N^*,
 \qquad
 \mathcal K_N=[B,AB,\ldots,A^{N-1}B].
\]
W szczególności
\(W_n^{\mathrm d}>0\)
dokładnie wtedy, gdy
para \((A,B)\) jest sterowalna.

Jeśli \(W_N^{\mathrm d}>0\),
to dla dowolnych
\(x_0,z\in\mathbb C^n\)
najmniejsza suma energii sterowania
\(\sum_{j=0}^{N-1}\|u_j\|_2^2\)
wśród ciągów
przeprowadzających
\(x_0\) w \(z\)
po \(N\) krokach
jest równa
\[
 (z-A^Nx_0)^*
 (W_N^{\mathrm d})^{-1}
 (z-A^Nx_0).
\]
Minimum jest osiągane
przez dokładnie jeden
ciąg sterowań.
\end{stwierdzenie}
\begin{proof}
Każdy składnik
\(A^jBB^*(A^*)^j\)
jest hermitowski
i ma nieujemną formę
kwadratową, gdyż
\[
 v^*A^jBB^*(A^*)^jv
 =\|B^*(A^*)^jv\|_2^2.
\]
Zatem ich suma
jest hermitowska
i nieujemnie określona.
Ponieważ
\(\mathcal K_N\)
jest złożona kolejno
z bloków \(A^jB\),
mnożenie przez
\(\mathcal K_N^*\)
daje dokładnie wskazaną sumę.
Dla dowolnej macierzy \(M\)
zachodzi
\(\ker(MM^*)=\ker M^*\):
jedna inkluzja jest oczywista,
a z \(MM^*v=0\) wynika
\(0=v^*MM^*v=\|M^*v\|_2^2\).
Stąd \(MM^*>0\)
dokładnie wtedy, gdy
\(M\) ma pełną rangę wierszową.
Dla \(N=n\) daje to
kryterium sterowalności.

Przejdźmy do kosztu.
Po odwróceniu kolejności
sterowań z
Twierdzenia~\ref{thm:lin25-kalman}
otrzymujemy
\[
 x_N=A^Nx_0+\mathcal K_Nw,
 \qquad
 w=\begin{pmatrix}
      u_{N-1}\\u_{N-2}\\\vdots\\u_0
    \end{pmatrix}.
\]
Niech \(d=z-A^Nx_0\).
Warunek osiągnięcia celu
ma postać \(\mathcal K_Nw=d\).
Jeśli \(W_N^{\mathrm d}>0\),
to definiujemy
\[
 w_0=\mathcal K_N^*
          (W_N^{\mathrm d})^{-1}d.
\]
Wtedy
\(\mathcal K_Nw_0
 =W_N^{\mathrm d}
    (W_N^{\mathrm d})^{-1}d=d\),
więc \(w_0\) jest dopuszczalny.
Każdy inny dopuszczalny
\(w\) ma postać \(w=w_0+h\)
z \(h\in\ker\mathcal K_N\).
Wektory \(w_0\) i \(h\)
są ortogonalne, gdyż
\[
 w_0^*h
 =d^*(W_N^{\mathrm d})^{-1}
    \mathcal K_Nh=0.
\]
Z twierdzenia Pitagorasa
wynika
\(\|w\|_2^2
 =\|w_0\|_2^2+\|h\|_2^2\),
a równość z wartością
minimalną zachodzi
tylko dla \(h=0\).
Wreszcie
\[
 \|w_0\|_2^2
 =d^*(W_N^{\mathrm d})^{-1}
    \mathcal K_N\mathcal K_N^*
    (W_N^{\mathrm d})^{-1}d
 =d^*(W_N^{\mathrm d})^{-1}d.
\]
To dowodzi wszystkich tez.
\end{proof}

\begin{przyklad}[Energia potrzebna do osiągnięcia celu]
\label{prz:lin25-energy}
Dla pary
\[
 A=\begin{pmatrix}0&1\\0&0\end{pmatrix},
 \qquad B=e_2
\]
mamy
\[
 W_2^{\mathrm d}
 =BB^*+ABB^*A^*
 =\begin{pmatrix}0&0\\0&1\end{pmatrix}
  +\begin{pmatrix}1&0\\0&0\end{pmatrix}
 =I_2.
\]
Jeśli \(x_0=0\)
i chcemy osiągnąć
\(z=(a,b)^{\mathsf T}\)
w dwóch krokach,
to optymalny ciąg sterowań
wynosi \(u_0=a\), \(u_1=b\),
a minimalna energia
\(|a|^2+|b|^2=\|z\|_2^2\).
Wynika to bezpośrednio
z ogólnego wzoru,
ale także z jawnej trajektorii
w Przykładzie~\ref{prz:lin25-chain}.
\end{przyklad}

\section{Przypadek rzeczywisty i granice twierdzeń}
\label{sec:lin25-real}

\begin{uwaga}[Macierze rzeczywiste]
\label{rem:lin25-real}
Definicje sterowalności
i obserwowalności
oraz kryteria rangowe
pozostają ważne dla macierzy
rzeczywistych.
W kryterium PBH
należy jednak testować
\emph{wszystkie zespolone}
wartości własne
rzeczywistej macierzy \(A\),
nie tylko rzeczywiste.

Wykażmy, że także
w twierdzeniu o stabilizowalności
można dla rzeczywistych \(A,B\)
wybrać \emph{rzeczywistą}
macierz \(K\).
W
Lemacie~\ref{lem:lin25-feedback-placement}
dla stabilizacji ciągłej
wystarczy wybierać
\(\mu_j\) z przedziału
\((-\infty,0)\) poza
skończonym widmem.
Dla dyskretnej
wystarczy przedział \((-1,1)\).
Argument w lemacie
używał zerowania funkcji
wymiernej na zbiorze
z nieskończenie wieloma punktami
skupienia, a więc działa
również na każdym z tych
rzeczywistych przedziałów.
W wybieranej kolejno
rzeczywistej bazie
można używać rzeczywistych
wektorów \(u_j\),
ponieważ macierze
\((H-\mu_jI)^{-1}D\)
są rzeczywiste.
Powstałe \(v_j\)
są wtedy rzeczywiste,
a odwzorowanie określone
przez \(Kv_j=u_j\)
ma rzeczywistą macierz.
W ten sposób
otrzymujemy rzeczywiste
sprzężenie stabilizujące,
bez wyprowadzania go
ze sprzężenia zespolonego.
Analogiczne rozumowanie
dla pary sprzężonej
daje rzeczywiste
wzmocnienie obserwatora \(L\).
\end{uwaga}

\begin{uwaga}[Przejście do ogólnych twierdzeń Riccatiego]
\label{rem:lin25-riccati}
Stabilizowalność nie oznacza
pełnej sterowalności:
nieosiągalne kierunki
mogą pozostać obecne,
o ile są stabilne.
Podobnie wykrywalność
dopuszcza stabilne kierunki
nieobserwowalne.
To właśnie te dwa
pojęcia występują
w ogólnych twierdzeniach
o stabilizujących rozwiązaniach
równań Riccatiego,
które rozszerzają
Rozdział~\ref{ch:lin24-riccati}
na niestabilną otwartą pętlę.
Ich szczegółowe sformułowania
i dowody pozostawiamy
kolejnemu rozdziałowi,
aby nie korzystać z nich
przed wprowadzeniem
koniecznych hipotez.
\end{uwaga}

\section{Podsumowanie}
\label{sec:lin25-summary}

Sterowalność określa,
czy sterowanie \(Bu\)
może wygenerować dowolny stan,
a obserwowalność, czy
obserwacja \(Cx\)
pozwala rozróżnić
wszystkie stany początkowe.
W czasie dyskretnym
charakteryzują je rangi
macierzy Kalmana
\[
 \mathcal K(A,B)=
 [B,AB,\ldots,A^{n-1}B],
 \qquad
 \mathcal O(C,A)=
 \begin{bmatrix}
 C\\CA\\\vdots\\CA^{n-1}
 \end{bmatrix}.
\]
W czasie ciągłym
te same kryteria wynikają
z szeregu wykładniczego;
dodatnia określoność
gramianu sterowalności
pozwala nawet podać
jawne sterowanie prowadzące
do dowolnego stanu.

Kryterium PBH zastępuje
duże macierze Kalmana
warunkiem rangi przy
wartościach własnych \(A\).
Rozkład Kalmana oddziela
kierunki osiągalne
od nieosiągalnych oraz
obserwowalne od nieobserwowalnych.
Macierz sprzężenia zwrotnego
może przesuwać wartości własne
części osiągalnej,
ale nie zmienia widma operatora
na ilorazie przez
przestrzeń osiągalną.
Dlatego stabilizowalność
wymaga dokładnie,
aby nieosiągalne wartości własne
były już stabilne.
Po zastosowaniu
związku dualnego
uzyskujemy analogiczne
kryterium wykrywalności
i zbieżność obserwatora.
